%==================================================================
% Paper P4 -- Section 2 -- Foundations (detailed, 5-7 pp)
%==================================================================
\section{Foundations}\label{sec:foundations}

This section reviews, in self-contained form and in the conventions
required by \Cref{thm:T1}, the chain of classical algebraicity and
integrality results on which the
Galois-Norm-Multiplicative-Invariant theorem rests:
\[
\text{Sch\"utt} \longrightarrow \text{Shimura} \longrightarrow
\text{Damerell} \longrightarrow \text{Hida} \longrightarrow
\text{Deligne} \longrightarrow \text{Kings-Sprang}.
\]
A reader familiar with this material may skim and proceed to
\Cref{sec:T1-proof}, after noting our specific conventions for
the period $\Omega_K$ in~\eqref{eq:chowla-selberg-2}, the
Tate-twist stratum~\eqref{eq:stratum-2}, and the canonical Hecke
character $\psi_K$ in~\Cref{ss:psi-K}.

\subsection{Sch\"utt's classification of $\Q$-rational CM newforms}
\label{ss:schutt}

Let $K = \Q(\sqrt{D})$ be imaginary quadratic with class group
$\Cl(K)$ a finite abelian $2$-group of $2$-rank $r$ (so $h_K =
2^r$). Sch\"utt's
classification~\cite{Schutt2008}, building on
Hecke~\cite{Hecke1937}, proves:

\begin{theorem}[Sch\"utt 2008; classical input]\label{thm:schutt}
The space $S_w(\Gamma_0(|D|), \chiK)^{\mathrm{CM},\Q\text{-rat}}$
of weight-$w$ CM newforms of level $|D|$ and Nebentypus $\chiK$
having $\Q$-rational Fourier coefficients satisfies
\begin{equation}\label{eq:schutt-dim-2}
\dim_\Q S_w(\Gamma_0(|D|), \chiK)^{\mathrm{CM},\Q\text{-rat}}
\;=\; |\Cl(K)[2]| \cdot \mathbf{1}_{\exp \Cl(K)\,\mid\, w-1}.
\end{equation}
The constructive identification realises the $\Q$-rational
eigenforms as theta series $\theta_\psi$ for algebraic Hecke
characters $\psi$ of $K$ of infinity type $(w-1, 0)$ whose
restriction to the class group is a \emph{genus character} (i.e.\
$\psi^2 = 1$ on $\Cl(K)$). The space of such genus characters has
dimension $|\Cl(K)[2]|$ by Pontryagin duality.
\end{theorem}

By Gauss's 1801 genus theorem~\cite{Gauss1801}~\cite[Thm.~3.15]{Cox1989},
\begin{equation}\label{eq:gauss-genus}
|\Cl(K)[2]| \;=\; 2^{t-1}, \qquad
t = \omega(|D|) = \#\{\text{prime divisors of }|D|\}.
\end{equation}
Throughout the remainder of the paper we work with $w \in \{3, 5\}$
and assume $\Cl(K) \cong (\Z/2)^r$, so both weight spaces have the
same dimension $h_K = 2^r$. For $r = 0$ (the nine $h_K = 1$
Heegner discriminants
$D \in \{-3, -4, -7, -8, -11, -19, -43, -67, -163\}$), the dimension
is~$1$, the cross-ratio theorem is vacuous, and the architecture
of~\Cref{thm:T1} must use the period normalisation directly, as
in~\Cref{thm:T2}. For $r = 1$ ($h_K = 2$), there are exactly two
normalised $\Q$-rational eigenforms in each weight, exchanged by
$\Gal(\HK/K) = \Cl(K) = \Z/2$.

\subsection{Canonical Hecke character and Galois orbits}
\label{ss:psi-K}

\begin{definition}[Canonical CM Hecke character]\label{def:psiK}
For $K = \Q(\sqrt{D})$ imaginary quadratic, fix an embedding $K
\hookrightarrow \C$ and choose a fundamental unit normalisation.
Define the canonical algebraic Hecke character $\psi_K \colon
\mathbb{A}_K^\times/K^\times \to \C^\times$ of infinity type
$(1, 0)$ by
\begin{equation}
\psi_K(\fr{a}) \;:=\; \alpha
\qquad \text{whenever } \fr{a} = (\alpha)\text{ is principal,
$\alpha \in K$, $\operatorname{Im}(\alpha) > 0$,}
\end{equation}
and extend uniquely to all integral ideals via class field theory
and ramification control at primes dividing $|D|$.
\end{definition}

For each integer $n \geq 1$, the $n$-th power $\psi_K^n$ is an
algebraic Hecke character of infinity type $(n, 0)$. For $h_K \geq
2$, the Galois action of $\sigma \in \Cl(K) \cong \Gal(\HK/K)$
gives the conjugate character
\begin{equation}
\psi_K^{\sigma}(\fr{a}) \;:=\; \psi_K(\sigma^{-1} \fr{a}),
\end{equation}
producing $h_K$ distinct Galois conjugates
$\psi_K^{(1)}, \dots, \psi_K^{(h_K)}$. The associated theta
series are
\begin{equation}\label{eq:theta-series}
f_w^{(j)}(\tau) \;:=\; \theta_{(\psi_K^{(j)})^{w-1}}(\tau)
\;=\; \sum_{\fr{a} \subset \OK} (\psi_K^{(j)}(\fr{a}))^{w-1} \cdot q^{\Nm(\fr{a})},
\qquad q := e^{2\pi i \tau},
\end{equation}
and by~\Cref{thm:schutt}, the $\Q$-rational ones (in our
$\Cl(K) = (\Z/2)^r$ setting) are exactly the $h_K$ Galois
conjugates $\{f_w^{(j)}\}$.

\subsection{Shimura algebraicity}\label{ss:shimura}

For a $\Q$-rational CM newform $f_w \in S_w(\Gamma_0(|D|), \chiK)$
attached to a Hecke character $\psi$ of $K$ of infinity type
$(w-1, 0)$, Shimura's algebraicity
theorem~\cite{Shimura1976} states:

\begin{theorem}[Shimura 1976]\label{thm:shimura}
For every critical integer $s \in \{1, 2, \dots, w-1\}$,
\begin{equation}\label{eq:shimura-2}
\frac{L(f_w, s) \cdot \Gamma(s)}{(2\pi i)^s \cdot \Omega_K^{w-1}}
\;\in\; \overline{\Q}.
\end{equation}
For $h_K = 1$, the algebraic part on the left lies in $\Q$
(Damerell~\cite{Damerell1971}). For $h_K \geq 2$, the algebraic
part lies in $\HK$
(Hida~\cite{Hida1981}; see~\Cref{ss:hida}).
\end{theorem}

\begin{remark}\label{rem:central-vs-critical}
Note that ``critical'' in Shimura's sense means $s \in
\{1, \dots, w-1\} \cap \Z$. The stratum
condition~\eqref{eq:stratum-2} below picks out the
\emph{Tate-balanced} pairs $(s_1, s_2)$ in the joint critical
strip for weights $w_1 = 5$, $w_2 = 3$, but is distinct from the
notion of ``central twist'' $s = w/2$, which is not an integer
for odd $w$. For odd $w \in \{3, 5\}$, the central twist lies on
the symmetry line of the functional equation but is not itself a
Shimura-critical point.
\end{remark}

\subsection{Damerell's explicit period and Chowla-Selberg formula}
\label{ss:damerell}

For $K$ imaginary quadratic of class number $h_K = 1$, the
Damerell period~\cite{Damerell1971} appearing in
\eqref{eq:shimura-2} admits the explicit Chowla-Selberg
description~\cite{ChowlaSelberg1949}
\begin{equation}\label{eq:chowla-selberg-2}
\Omega_K \;=\; \sqrt{\pi} \cdot |D|^{-1/4} \cdot
\prod_{a=1}^{|D|-1} \Gamma\!\bigl(\tfrac{a}{|D|}\bigr)^{\chiK(a)/(2h_K)}.
\end{equation}
For $h_K \geq 2$, the period $\Omega_K$ in~\eqref{eq:shimura-2}
must be replaced by an $\HK$-valued period that depends on the
specific CM elliptic curve in the relevant Galois orbit; see
\cite[\S 8]{Hida1981}.

\begin{example}[Numerical verification at the seven non-vanishing
Heegner anchors]\label{ex:omega-K-numerics}
Direct PARI/GP evaluation of~\eqref{eq:chowla-selberg-2} at $80$-digit
precision (\texttt{realprecision = 80}) for the seven non-vanishing
Heegner discriminants gives (abbreviated to $14$~digits):
\[
\begin{array}{|r|l|l|}
\hline
D & \Omega_K & \Omega_K \cdot |D|^{1/4} \\
\hline
-7   & 3.61689501212884 & 5.88275 \\
-8   & 3.56230079222662 & 5.99005 \\
-11  & 3.39582468290106 & 6.18476 \\
-19  & 2.96316633590618 & 6.18816 \\
-43  & 2.04392119260363 & 5.23218 \\
-67  & 1.49116290156061 & 4.27027 \\
-163 & 0.56119333550365 & 2.00428 \\
\hline
\end{array}
\]
The quantity $\Omega_K \cdot |D|^{1/4}$ is bounded in
$[2.0, 6.2]$ across the seven anchors, consistent with the
Stirling-type asymptotic $\Omega_K = O(|D|^{-1/4})$ predicted
by~\eqref{eq:chowla-selberg-2}.
\end{example}

\subsection{Hida's Hilbert-class-field rationality at $h_K \geq 2$}
\label{ss:hida}

For $h_K \geq 2$, Hida~\cite{Hida1981} proves the natural
$\HK$-equivariant extension of Shimura's algebraicity:

\begin{theorem}[Hida 1981]\label{thm:hida}
For $f_w$ a $\Q$-rational CM newform attached to a Galois orbit
of Hecke characters $\{\psi_K^{(j)}\}_{j=1}^{h_K}$ on $K$ of
infinity type $(w-1, 0)$, the algebraic part
\[
A_w^{(j)}(s) \;:=\; \frac{L(f_w^{(j)}, s) \cdot \Gamma(s)}{(2\pi i)^s \cdot \Omega_K^{w-1}}
\]
lies in $\HK$ for every critical $s$, and the Galois action of
$\sigma \in \Gal(\HK/K) \cong \Cl(K)$ on $A_w^{(j)}(s)$ is by
permutation of the Galois orbit:
\begin{equation}\label{eq:hida-galois-action}
\sigma\bigl(A_w^{(j)}(s)\bigr) \;=\; A_w^{(\sigma \cdot j)}(s).
\end{equation}
\end{theorem}

\begin{corollary}[$K$-rationality of the multiplicative norm]
\label{cor:K-rational-norm}
The norm
$\Nm_{\HK/K}\bigl(A_w^{(1)}(s)\bigr) := \prod_{\sigma \in \Cl(K)}
A_w^{(\sigma \cdot 1)}(s)
\in K$. Since $f_w^{(j)}$ has real (in fact integral)
Fourier coefficients ($\Q$-rationality), complex conjugation acts
trivially on $A_w^{(j)}(s)$ for real critical $s$, hence
$\Nm_{\HK/K}\bigl(A_w^{(1)}(s)\bigr) \in K \cap \R = \Q$ (as $K
= \Q(\sqrt{D})$ with $D < 0$).
\end{corollary}

\Cref{cor:K-rational-norm} is the technical engine behind
Step~3 (Galois descent) and Step~4 (reality) of the proof
of~\Cref{thm:T1} in~\Cref{ss:proof-sketch} below.

\subsection{Deligne's Tate-twist exponent calculus}
\label{ss:deligne-tate}

Deligne's 1979 conjecture~\cite{Deligne1979}, classical in our
setting since the motives involved are $1$-dimensional Hecke
characters, attaches to each motive $M$ over $\Q$ critical at $s$
a canonical period $c^\pm(M)$ encoding the Hodge-de~Rham
comparison. The Tate-twist functor $M \mapsto M(j)$ satisfies the
multiplicative relation
\begin{equation}\label{eq:tate-twist}
c^\pm(M(j)) \;=\; (2\pi i)^j \cdot c^\pm(M),
\qquad \Q(j_1) \otimes \Q(j_2) \;=\; \Q(j_1 + j_2).
\end{equation}
For our two CM motives $M_i := M(f_{w_i})$ of Hodge weight
$w_i - 1$ (Hodge type $(w_i - 1, 0) \oplus (0, w_i - 1)$), a
product of $L$-values $L(M_1, s_1) \cdot L(M_2, s_2)$ has period
factor $(2\pi i)^{s_1 + s_2}$ times $\Omega_K^{(w_1 - 1) + (w_2 - 1)}$.
In a cross-ratio of the form
\begin{equation}\label{eq:cross-ratio-form}
\frac{L(M_1^{(1)}, s_1) \cdot L(M_2^{(2)}, s_2)}
     {L(M_1^{(2)}, s_1) \cdot L(M_2^{(1)}, s_2)},
\end{equation}
both Tate-twist exponents and both Damerell period exponents
\emph{cancel exactly}, provided one chooses two pairs of equal
$s_1 + s_2$ and equal $(w_1 - 1) + (w_2 - 1)$. For weights
$\{w_1, w_2\} = \{3, 5\}$, the Hodge weight condition forces
$(w_1 - 1) + (w_2 - 1) = 6$, automatic; the Tate condition forces
the \emph{stratum condition}
\begin{equation}\label{eq:stratum-2}
s_1 + s_2 \;=\; \frac{w_1 + w_2}{2} \;=\; 4.
\end{equation}
This is the structural origin of the stratum
condition~\eqref{eq:stratum} in~\Cref{thm:T1}: without it, the
Tate-twist factors do not balance and the ratio is no longer
period-clean.

\subsection{Kings-Sprang integrality}\label{ss:kings-sprang}

The very recent theorem of Kings and Sprang~\cite{KingsSprang2025}
(\cite{KingsSprangArxiv}), building on Bannai-Kobayashi's
Eisenstein-Kronecker theta classes on the Poincar\'e
bundle~\cite{BannaiKobayashi2010}, sharpens
Damerell-Shimura algebraicity to \emph{integrality}:

\begin{theorem}[Kings-Sprang 2025]\label{thm:kings-sprang}
Let $\psi$ be an algebraic Hecke character of $K$ of infinity
type $(w, 0)$ with $w \geq 1$, and let $s_0$ be a critical
integer for $L(\psi, s)$. Then
\begin{equation}\label{eq:kings-sprang}
\frac{L(\psi, s_0) \cdot \Gamma(s_0)}{(2\pi i)^{s_0} \cdot \Omega_K^{w}}
\;\in\; \mathcal{O}_{\HK},
\end{equation}
modulo torsion factors controlled by the explicit Tate twist. The
denominator of the integral relation is bounded by the
Eisenstein-Kronecker denominator, supported only at ramification
primes of $K$ and at small primes appearing in the conductor of
$\psi$.
\end{theorem}

\Cref{thm:kings-sprang} is the quantitative version of
\eqref{eq:shimura-2} that yields the denominator bound
$|D|^{2d}$ in~\Cref{thm:T1}. Indeed, applying
\Cref{thm:kings-sprang} to each of the $d$ factors of $L(f_{w_1}^{(j)},
s_1)$ and $d$ factors of $L(f_{w_2}^{(j)}, s_2)$ in the polynomial
$P$ of bidegree $(d, d)$, the resulting rational evaluation has
denominator dividing
$\Nm_{K/\Q}(\disc \OK)^d = |D|^{2d}$, up to the explicitly
bounded local Eisenstein-Kronecker correction.

\subsection{Newton-Dickson cross-weight chain}
\label{ss:newton-dickson}

The Hecke recursion~\cite{Hecke1937}~\cite{Schutt2008}
interpolating weight-$3$, weight-$5$, weight-$7$, \ldots\
eigenforms in a Hida family~\cite{Hida1986} reads, for split
primes $p$ of $K$:

\begin{proposition}[Newton-Dickson chain]\label{prop:newton-dickson}
For $f_{2m+1} = \theta_{\psi_K^{2m}}$ the weight-$(2m+1)$ CM
newform attached to the $2m$-th power of the canonical Hecke
character of \Cref{def:psiK}, and $p$ a rational prime split in
$K$,
\begin{equation}\label{eq:newton-dickson}
a_p(f_{2m+1}) \;=\; D_m\bigl(a_p(f_3), p\bigr),
\end{equation}
where $D_m$ is the $m$-th modified Dickson polynomial defined by
\[
D_1(x, p) = x, \quad D_2(x, p) = x^2 - 2p^2,
\quad D_m(x, p) = x \cdot D_{m-1}(x, p) - p^2 D_{m-2}(x, p)
\quad (m \geq 3).
\]
\end{proposition}

\begin{remark}[Numerical verification]
At $D = -67$ (a $h_K = 1$ Heegner anchor), four split primes
$p \in \{17, 19, 23, 29\}$ give:
\[
\begin{array}{|r|r|r|r|}
\hline
p & a_p(f_3) & a_p(f_5) & a_p(f_7) \\
\hline
17 & -33 & 511 & -7\,326 \\
19 & -29 & 119 & 7\,018 \\
23 & -21 & -617 & 24\,066 \\
29 & -9 & -1\,601 & 21\,978 \\
\hline
\end{array}
\]
Each row satisfies $a_p(f_5) = a_p(f_3)^2 - 2p^2$ and
$a_p(f_7) = a_p(f_3)^3 - 3 p^2 \cdot a_p(f_3)$ exactly, confirming
$D_2, D_3$ at $4 + 4 = 8$ data points. Cumulative verification
across six Heegner discriminants $D \in \{-7, -11, -19, -43, -67,
-163\}$, weights $w \in \{3, 5, 7\}$, and all split primes
$p \leq 30$ gives $81/81$ exact matches at $80$-digit PARI
precision.
\end{remark}

\subsection{Bridge B$'$: functional-equation identity}
\label{ss:bridge-Bprime}

A useful sanity check, verified independently at $96$-digit PARI
precision across all seven non-vanishing Heegner anchors plus
$D = -8$, is the cross-weight functional-equation identity

\begin{proposition}[Bridge B$'$, FE identity]\label{prop:bridge-Bprime}
For every $\Q$-rational weight-$3$ CM newform $f_3$ and the
corresponding weight-$5$ newform $f_5$ attached to the same
$K = \Q(\sqrt{D})$ (so $f_5 = \theta_{\psi_K^4}$, $f_3 = \theta_{\psi_K^2}$),
\begin{equation}\label{eq:bridge-Bprime}
\frac{L(f_5, 3)}{L(f_3, 1)} \;=\; \frac{2\pi^2}{|D|} \cdot
\frac{L(f_5, 2)}{L(f_3, 2)}.
\end{equation}
\end{proposition}

\begin{proof}[Sketch]
Apply the Hecke completed functional equation
$\Lambda(f_w, s) = \varepsilon \cdot \Lambda(f_w, w - s)$,
with $\Lambda(f_w, s) := |D|^{s/2} \cdot (2\pi)^{-s} \cdot
\Gamma(s) \cdot L(f_w, s)$, separately to $f_3$ at $s_3 = 1
\leftrightarrow 2$ and to $f_5$ at $s_5 = 3 \leftrightarrow 2$,
and take the ratio. The $|D|$ and $(2\pi)$ powers combine to
give $2\pi^2/|D|$.
\end{proof}

\Cref{prop:bridge-Bprime} is the key bridge used in the
unconditional lower-bound proof of~\Cref{thm:T2}; it transports
the Euler-product bound at $s = 3$ (where convergence is
manifest, $3 > w/2 + 1/2 = 2$) to a bound at $s = 1$
(within the critical strip).

\subsection{Synopsis: the period diagram}\label{ss:diagram}

We collect the six ingredients into a single commutative diagram
describing the period quotient that underlies~\Cref{thm:T1}.

\begin{equation}\label{diag:period-2}
\begin{tikzcd}[row sep=large, column sep=large]
L(f_w^{(j)}, s)
\arrow[r, "{\;\Gamma(s)/((2\pi i)^s \Omega_K^{w-1})\;}"]
\arrow[d, "{\text{Kings-Sprang}}"']
& A_w^{(j)}(s) \in \HK
\arrow[d, "{\text{Hida}}"] \\
\mathcal{O}_{\HK}
\arrow[r, "{\Nm_{\HK/K}}"']
& \HK^{\Cl(K)} = K \\
\end{tikzcd}
\end{equation}

\medskip
\noindent\emph{Reading.} The top row is Shimura algebraicity:
each $L$-value, after dividing by the universal
$(2\pi i)^s \Omega_K^{w-1}$ factor, lands in an algebraic
number, refined by Hida to live in $\HK$ for $h_K \geq 2$. The
left vertical arrow is Kings-Sprang: the algebraic part is in
fact an algebraic integer of $\HK$ (up to bounded denominator).
The bottom row is Galois descent: the multiplicative norm of an
$\HK$-element down to $K$ lands in the Galois-invariants
$\HK^{\Cl(K)} = K$. For Galois-equivariant polynomials in
$h_K$ such elements, the norm is replaced by an arbitrary
$\Cl(K)$-symmetric polynomial, with the same conclusion.

The right-hand column ends in $\HK$ for individual values;
Galois descent collapses to $K$ via the diagonal
$\Cl(K)$-action, and reality plus the imaginary-quadratic
hypothesis collapses $K$ further to $\Q$. This is the abstract
shape of the argument made precise in~\Cref{ss:proof-sketch}.

\subsection{Tate-Deligne motives of Hecke characters}
\label{ss:motives}

For later reference (and for the statement of \Cref{thm:T3-intro}
on the partial functor), we recall the motivic upgrade of the
analytic objects above. By Tate-Deligne 1965 (classical;
Bourbaki S\'eminaire 306), for every algebraic Hecke character
$\psi \colon \mathbb{A}_K^\times/K^\times \to \C^\times$ of
infinity type $(a, b)$ with $a + b \geq 0$, there exists a
$1$-dimensional rational motive $M(\psi) \in \DM_\Q$ of weight
$a + b$ whose $L$-function is $L(\psi, s)$ and whose realisations
are:
\begin{itemize}
\item \emph{$\ell$-adic.} The $1$-dimensional $\ell$-adic Galois
representation $\Gal(\overline{K}/K) \to \overline{\Q_\ell}^\times$
attached to $\psi$, extended by induction to
$\Gal(\overline{\Q}/\Q) \to \mathrm{GL}_2(\overline{\Q_\ell})$.
\item \emph{Hodge.} The $1$-dimensional Hodge structure of weight
$a + b$ with Hodge filtration concentrated in $(a, b)$.
\item \emph{de Rham.} A $1$-dimensional $\Q$-vector space with
the standard comparison isomorphism to Hodge.
\end{itemize}
The Tannakian formalism gives $M(\psi) \otimes M(\psi') = M(\psi
\cdot \psi')$. For our setting $\psi = \psi_K^n$, the motive
$M(\psi_K^n)$ is the Scholl~\cite{Scholl1990} motive attached
to the CM theta series $\theta_{\psi_K^n}$ of weight $n + 1$:
\begin{equation}\label{eq:scholl-CM}
M(\psi_K^n) \;\cong\; M(\theta_{\psi_K^n}) \;\in\; \DM_\Q,
\qquad n \geq 1.
\end{equation}
This motive carries a Deligne polarisation $c^+(M(\psi_K^n))
\in \C^\times / \Q^\times$, identified with $\Omega_K^n$ under
the conventions of \Cref{ss:damerell}. The Tate twist
$M(\psi_K^n)(j)$ has polarisation $(2\pi i)^j \cdot \Omega_K^n$,
reflecting~\eqref{eq:tate-twist}.

We package the foundations into the following statement, which
underlies both the cross-ratio theorem of~\Cref{sec:T1-proof} and
the partial functor of~\Cref{sec:T3}:

\begin{proposition}[Period-clean cross-ratio criterion]
\label{prop:period-clean}
Let $\psi_1, \psi_2$ be algebraic Hecke characters of $K$ of
infinity types $(w_1 - 1, 0)$, $(w_2 - 1, 0)$, and let
$s_1, s_2$ be critical integers for $L(\psi_i, s)$ satisfying
the stratum condition $s_1 + s_2 = (w_1 + w_2)/2$. For each
$\sigma \in \Cl(K)$, the multiplicative cross-ratio
\[
\frac{L(\psi_1^{(1)}, s_1) \cdot L(\psi_2^{(2)}, s_2)}
     {L(\psi_1^{(2)}, s_1) \cdot L(\psi_2^{(1)}, s_2)}
\;=\;
\frac{A_1^{(1)}(s_1) \cdot A_2^{(2)}(s_2)}
     {A_1^{(2)}(s_1) \cdot A_2^{(1)}(s_2)}
\;\in\;\HK,
\]
with the $(2\pi i)^{s_1 + s_2} \cdot \Omega_K^{(w_1 - 1) + (w_2 - 1)}$
period factor cancelling in numerator and denominator.
\end{proposition}

\Cref{prop:period-clean} is the central simplification that makes
the proof of \Cref{thm:T1} possible in elementary steps; in the
absence of the stratum condition, no such period-clean cross-ratio
exists, and the conclusion of \Cref{thm:T1} fails (as we verify
empirically in~\Cref{ss:no-closed-form}).

\medskip
\noindent
With the foundations in place we proceed to the proof
of~\Cref{thm:T1} (\Cref{sec:T1-proof}).
